\documentclass[11pt]{article}
\usepackage[margin=1in]{geometry}
\usepackage[T1]{fontenc}
\usepackage[utf8]{inputenc}
\usepackage{lmodern}
\usepackage{microtype}
\usepackage{enumitem}
\usepackage{booktabs}
\usepackage{array}
\usepackage{amsmath}
\usepackage{amssymb}
\usepackage{hyperref}
\usepackage{natbib}

\title{Sparse Features for Surface-Cue Reliance in Gemma-Based Reward Judges}
\author{AISafety project draft}
\date{\today}

\begin{document}
\maketitle

\begin{abstract}
LLM-based judges often favor LLM-generated text over human text in realistic
evaluation settings. This draft reports an ongoing mechanistic study of one
possible substrate for that behavior: sparse residual-stream features that
encode surface and discourse cues such as structured assistant formatting,
formal institutional phrasing, and benefit/value framing. We construct a
linguistically grounded atom-to-bundle ontology, discover candidate Gemma Scope
SAE features for those atoms, test whether the frozen candidates align with a
Gemma 2 reward judge's human-vs-LLM reward margins across HC3, HC3+, H-LLMC2,
and HAP-E, and then mute selected feature bundles at inference time. The
current strongest result is a causal perturbation of the
\emph{structured assistant packaging} bundle: muting ten source-stable SAE
features lowers the judge's LLM-minus-human reward margin by 0.0273 relative
to matched random feature controls, with a 0.0410 shift in high-occurrence
pairs and a 2.4 percentage point reduction in thresholded LLM preference.
Formal/institutional packaging shows a smaller but specific margin effect,
while benefit/value framing is weak and control-confounded. We treat these
results as evidence for feature-level surface-cue reliance, not yet as a full
repair method; broad retention and combined-intervention tests remain the next
claim gates.
\end{abstract}

\section{Introduction}

LLM judges and reward models are increasingly used as scalable proxies for
human preference judgments \citep{zheng2023mtbench}. This scaling benefit comes
with a familiar problem: judges can respond to cues that correlate with
authorship or presentation rather than with task quality. Prior work documents
position bias, verbosity bias, self-preference, and AI--AI bias in which LLM
judges favor LLM-generated communications \citep{laurito2025aiaibias,
panickssery2024selfrecognition,lee2024remodetect}. These findings establish
that evaluator bias is real, but they do not identify the internal features
that carry the bias.

This project asks whether part of the LLM-favoring margin in a Gemma-based
reward judge is mediated by sparse features for surface-cue bundles. The goal
is not to show that every use of formatting or style is invalid. Structured
answers can genuinely improve readability. The narrower question is whether
specific sparse features for presentation form measurably affect reward margins
on content-matched human-vs-LLM pairs, and whether those effects are specific
relative to matched random feature perturbations.

The current draft focuses on a single anchor judge, denoted J0, built from
Gemma 2 9B \citep{gemmateam2024gemma2}. We use Gemma Scope residual-stream
SAEs as the feature basis \citep{lieberum2024gemmascope}. The study proceeds in
four stages:

\begin{enumerate}[leftmargin=1.5em]
\item build a linguistic atom-to-bundle ontology of surface cues;
\item discover and freeze SAE features associated with those atoms and bundles;
\item test whether frozen candidates align with J0 human-vs-LLM reward margins;
\item mute selected feature bundles and compare reward-margin shifts against
      layer-matched random controls.
\end{enumerate}

The findings are preliminary but structured enough to discuss. The strongest
causal result is not broad ``formality'' but a bundle we call
\emph{structured assistant packaging}: bullet/list structure, multi-part
completion, helpful wrapper phrasing, sentence-length balance, and
compression/elaboration. Muting this bundle produces a much larger and more
source-stable reward-margin shift than muting formal/institutional packaging or
benefit/value framing.

\section{Related Work}

\paragraph{LLM judges and AI--AI bias.}
LLM-as-a-judge systems can approximate preference judgments at scale, but their
failure modes include position effects, verbosity effects, self-preference, and
other systematic biases \citep{zheng2023mtbench,panickssery2024selfrecognition}.
The motivating ecological result is AI--AI bias: LLM judges favor
LLM-generated paper abstracts, product descriptions, and movie synopses over
human alternatives \citep{laurito2025aiaibias}. Reward models can also
recognize aligned-model generations \citep{lee2024remodetect}, suggesting that
authorship-correlated signals can be available internally to evaluators.

\paragraph{Linguistic grounding.}
Our ontology is motivated by register theory, computational sociolinguistics,
and metadiscourse rather than by ad hoc prompt labels
\citep{biber1993register,nguyen2016sociolinguistics,passonneau2014biberredux,
hyland2005metadiscourse}. Recent human-vs-LLM corpora also show that rhetorical
and grammatical features remain discriminative in modern generated text
\citep{reinhart2024llmshumans}. We use HC3 \citep{guo2023hc3}, HC3+, H-LLMC2
\citep{noepsl2025hllmc2}, and HAP-E \citep{browndw2024hape,
reinhart2024llmshumans} to avoid defining the ontology solely from Laurito-style
examples.

\paragraph{Sparse feature methods.}
Sparse autoencoders expose interpretable feature bases in transformer models
\citep{cunningham2023sae,gao2024scaling}. Gemma Scope provides a public SAE
basis for Gemma 2 \citep{lieberum2024gemmascope}. Recent work on sparse feature
circuits and SAE-based model editing motivates using feature-level interventions
as causal tests rather than stopping at feature correlation
\citep{marks2024sparsecircuits,dunefsky2024transcoders,kissane2024attention}.
This paper uses SAE feature muting as a first causal readout, not as a complete
circuit analysis.

\section{Experimental Setup}

\subsection{Judge And Feature Basis}

The anchor judge J0 is a Gemma 2 9B based scalar reward model with a LoRA
adapter and value head. A text receives a scalar reward from the final
non-padding token representation. A human-vs-LLM pair has margin
\[
  m = r(x_{\mathrm{LLM}}) - r(x_{\mathrm{human}}).
\]
Positive margin means the reward model favors the LLM text. We report
continuous margins as the primary metric. The thresholded ``LLM chosen'' rate,
\(\mathbb{1}[m > 0]\), is included only as a secondary diagnostic because it is
an arbitrary scalar-model decision boundary rather than an independently
elicited binary judge choice.

For feature interpretation and perturbation we use Gemma Scope residual SAEs.
Hidden-state index \(h\) maps to SAE layer \(h-1\). Candidate features are
aggregated over valid tokens by max activation, matching the discovery and
alignment pipeline.

\subsection{Atom-to-Bundle Ontology}

The primitive unit is an \emph{atom}: a local, linguistically interpretable cue
such as bullet/list structure, formal connectives, or promotional adjectives.
Atoms are grouped into broader \emph{bundles} only when multiple atoms form a
coherent presentation family. The current frozen intervention registry contains
three primary bundles:

\begin{table}[h]
\centering
\small
\begin{tabular}{p{0.25\linewidth}p{0.48\linewidth}r}
\toprule
Bundle & Atoms & Target features \\
\midrule
Structured assistant packaging &
bullet/list structure; balanced multi-part completion; helpful assistant
wrapper; sentence-length balance; compression/elaboration &
10 \\
Formal/institutional packaging &
formal connectives; institutional impersonality; professional
self-presentation; nominalization patterns; complex noun phrase chains &
9 clean / 10 with source-sensitive add-back \\
Benefit/value framing &
benefit-first packaging; feature-benefit scripts; promotional adjectives;
enthusiasm/salesmanship &
9 \\
\bottomrule
\end{tabular}
\caption{Frozen bundle families used for current perturbation experiments.}
\label{tab:bundles}
\end{table}

\subsection{Broad Human-vs-LLM Alignment Set}

The broad confirmation set contains 10,000 content-matched human-vs-LLM pairs:
2,797 HC3 pairs, 2,500 HC3+ pairs, 1,374 H-LLMC2 pairs, and 3,329 HAP-E pairs.
J0 favors the LLM text in 78.57\% of these pairs, with mean LLM-minus-human
margin 0.3317. This set is not used to manually craft perturbations; it
provides the common evaluation substrate for candidate alignment and feature
muting.

\subsection{Candidate Discovery And Freeze}

The discovery stage combines atom-level SAE localization, bundle-level
aggregation, Laurito decision alignment, and broad human-vs-LLM alignment. The
frozen registry contains 29 bundle-feature rows, 75 atom-feature rows, and 87
matched random-control rows. It marks features as intervention-eligible when
they exceed matched random-control baselines and show adequate source support.
The registry also marks one formal feature, L34F7691, as source-sensitive.

\begin{table}[h]
\centering
\small
\begin{tabular}{lrrrrr}
\toprule
Bundle & Features & Eligible & Source-sensitive & Max \(|\rho_c|\) & Mean \(|\rho_c|\) \\
\midrule
Structured assistant packaging & 10 & 10 & 0 & 0.451 & 0.296 \\
Formal/institutional packaging & 10 & 9 & 1 & 0.437 & 0.220 \\
Benefit/value framing & 9 & 9 & 0 & 0.350 & 0.197 \\
\bottomrule
\end{tabular}
\caption{Frozen registry summary. \(\rho_c\) is the length/source-controlled
Spearman correlation between feature activation deltas and J0 reward-margin
deltas in the broad alignment run.}
\label{tab:registry}
\end{table}

\subsection{Feature Muting Protocol}

For a target bundle, we load the bundle's SAE features and register forward
hooks at the corresponding decoder layers. For each token representation
\(z\), the hook encodes \(z\) with the SAE, selects target feature activations,
and subtracts the selected decoder-vector contributions. The 1.0 mute used in
the current main runs is:
\[
  z' = z - \sum_{f \in F} a_f(z) W^{\mathrm{dec}}_f.
\]
For every target run, we also mute layer-matched random SAE features from the
frozen registry. The main reported effect is:
\[
  \Delta_{\mathrm{target-control}}
  =
  \left(m_{\mathrm{target\ mute}} - m_{\mathrm{baseline}}\right)
  -
  \left(m_{\mathrm{random\ mute}} - m_{\mathrm{baseline}}\right).
\]
Negative values mean the target mute lowers the LLM-favoring margin more than
the matched random mute.

Each perturbation run uses the same 1,200-pair subset of the broad alignment
set, with 165 H-LLMC2, 404 HAP-E, 337 HC3, and 294 HC3+ pairs. Pairs are binned
by the target bundle's signed activation-delta score into high, middle, and low
occurrence bins (300/600/300).

\section{Results}

\subsection{Structured Assistant Packaging Is The Strongest Causal Bundle}

Muting structured assistant packaging produces a large, specific, and
source-stable reduction in J0's LLM-minus-human margin. The all-pair
target-control shift is -0.0273. In high-occurrence pairs the shift is -0.0410,
while matched random controls are essentially flat.

\begin{table}[h]
\centering
\small
\begin{tabular}{lrrrr}
\toprule
Run & All pairs & High bin & Low bin & LLM-choice shift \\
\midrule
Structured assistant packaging & -0.0273 & -0.0410 & -0.0088 & -2.42 pp \\
Formal/institutional, clean & -0.0048 & -0.0087 & -0.0011 & -0.25 pp \\
Formal + L34F7691 add-back & -0.0038 & -0.0140 & +0.0078 & +0.42 pp \\
Benefit/value framing & -0.0028 & -0.0032 & +0.0012 & -1.83 pp \\
\bottomrule
\end{tabular}
\caption{1.0 feature mute effects. Entries are target-control changes in
LLM-minus-human reward margin, except the last column which reports the target
mute's thresholded LLM-choice-rate change.}
\label{tab:perturbation-main}
\end{table}

The structured effect is substantially larger than the other bundles. On the
same pair set, the structured all-pair target-control shift is about 5.6 times
the clean formal shift and 9.7 times the benefit/value shift. The effect is
also occurrence-sensitive: the high-minus-low target-control contrast is
-0.0322, and the Spearman correlation between signed bundle occurrence and the
target-control margin shift is -0.441.

\subsection{The Structured Effect Is Source-Stable}

Structured assistant packaging moves all four source families in the expected
direction:

\begin{table}[h]
\centering
\small
\begin{tabular}{lrrrr}
\toprule
Bundle & H-LLMC2 & HAP-E & HC3 & HC3+ \\
\midrule
Structured assistant packaging & -0.0421 & -0.0273 & -0.0309 & -0.0149 \\
Formal/institutional, clean & -0.0113 & -0.0018 & -0.0055 & -0.0046 \\
Formal + L34F7691 add-back & -0.0159 & +0.0037 & -0.0058 & -0.0051 \\
Benefit/value framing & -0.0012 & -0.0110 & +0.0014 & +0.0028 \\
\bottomrule
\end{tabular}
\caption{Source-specific target-control margin shifts. Negative values indicate
that the target mute lowers the LLM-favoring margin more than matched random
controls.}
\label{tab:source-effects}
\end{table}

This source table is important for interpretation. Structured packaging remains
stable even on HAP-E, while formal packaging is weaker there and benefit/value
framing is unstable across sources.

\subsection{Formal Packaging Is Specific But Modest}

The clean formal/institutional bundle has a small but directionally coherent
effect: -0.0048 all-pair target-control shift and -0.0087 in high-occurrence
pairs. Matched random controls remain near zero. This supports a secondary
claim that formal packaging features contribute to reward margins, but the
effect is too small to carry the paper alone.

The source-sensitive add-back run confirms the registry warning. Adding
L34F7691 increases the high-bin formal effect from -0.0087 to -0.0140, but it
also reverses the low bin to +0.0078 and flips HAP-E positive (+0.0037).
Therefore L34F7691 is useful as an ablation but should not be part of the clean
main formal result.

\subsection{Benefit/Value Framing Is Weak And Control-Confounded}

Benefit/value framing is not currently a strong causal bundle. The target mute
lowers margins by -0.0146, but matched random controls also lower margins by
-0.0118, leaving only -0.0028 target-control effect. The source pattern is
unstable: HAP-E moves negative, but HC3 and HC3+ move positive after
target-control correction. This bundle is best treated as a weak contrast or
as evidence that correlation-based feature selection does not guarantee a
specific causal effect.

\section{Interpretation}

The current results support a more specific story than ``LLM judges like
formal writing.'' The strongest causal evidence points to structured assistant
packaging: a package of list structure, multi-part answer formatting, helpful
wrapper phrasing, and balanced elaboration. This is a surface-cue bundle, but
it is not necessarily meaningless. Structured formatting can improve perceived
readability and may sometimes improve actual usefulness. The defensible claim
is that sparse features for this presentation bundle materially affect the
reward judge's LLM-favoring margin on content-matched human-vs-LLM pairs.

The formal/institutional bundle provides a useful secondary result because it
is less reducible to list formatting. Its effect is smaller but specific. The
source-sensitive add-back is also informative: a high-alignment feature can
increase high-bin sensitivity while damaging source and low-bin stability. This
validates the registry's source-sensitivity flag and argues for keeping
intervention candidates frozen and auditable.

The benefit/value result is a negative or weak result. It suggests that some
features discovered through atom alignment and broad margin correlation are not
good intervention targets. A future selection rule should probably incorporate
causal effect size directly, rather than ranking only by atom detection and
margin correlation.

\section{What The Current Evidence Does And Does Not Show}

The current evidence supports:

\begin{itemize}[leftmargin=1.5em]
\item SAE features can be organized into interpretable atom-to-bundle
      hypotheses.
\item Frozen structured assistant packaging features are strongly aligned with
      J0 human-vs-LLM margins.
\item Muting those features selectively reduces LLM-favoring reward margins
      relative to matched random feature controls.
\item The structured effect is present across HC3, HC3+, H-LLMC2, and HAP-E.
\item Formal/institutional features show a smaller but specific effect.
\end{itemize}

The current evidence does not yet show:

\begin{itemize}[leftmargin=1.5em]
\item that feature muting preserves general preference ability;
\item that the structured features are never utility-relevant;
\item that this is a complete circuit for AI--AI bias;
\item that benefit/value framing is a stable causal bundle;
\item that a combined intervention is an acceptable repair.
\end{itemize}

\section{Next Experiments}

The next required experiment shifts from SAE features as the primary actuator
to SAE features as interpretable detectors. First, deterministic
surface-counterfactual rewrites should increase or decrease one presentation
axis while approximately preserving content. The audit then measures
\(R(x_{\mathrm{variant}})-R(x_{\mathrm{base}})\), length-controlled reward
deltas, source consistency, and whether the matching SAE bundle activation
moves in the expected direction. This directly asks whether J0 rewards
structured assistant packaging, formal/institutional packaging, or benefit
framing when the substantive answer is held approximately fixed.

Second, the reward readout should be tested with pooled-state nulling. For each
axis, we fit a direction from counterfactual hidden-state deltas and project
the resulting subspace out before J0's scalar value head. This stronger
intervention should report three quantities: reduction in counterfactual
surface-cue reward deltas, change in broad human-vs-LLM reward margins, and
retention on ordinary preference pairs. SAE feature muting remains a useful
diagnostic baseline, but readout nulling is the more appropriate mitigation
test if individual sparse latents are too weak or redundant as actuators.

The paper should only move from diagnosis to mitigation language if readout
nulling or a later surface-invariant adapter reduces counterfactual cue
sensitivity while preserving most general preference behavior.

\section{Limitations}

The study currently uses one Gemma 2 9B anchor judge and one public SAE basis.
The feature-muting method subtracts SAE decoder-vector contributions and is not
a full circuit intervention. Max-pooled feature activation is useful for bundle
occurrence but may miss token-position specificity. Deterministic
counterfactuals are auditable but not perfect semantic controls, and
readout-space nulling is less mechanistically localized than SAE feature
damping. The human-vs-LLM pair set is broad but not a substitute for
task-specific human preference labels. Most importantly, reduced LLM-favoring
margin is not automatically improved judgment: the intervention must pass
retention tests before being described as repair.

\section{Provisional Paper Narrative}

A workshop paper can make the following compact claim:

\begin{quote}
We build a linguistic atom-to-bundle ontology for surface cues in
human-vs-LLM evaluation, localize candidate cues in Gemma Scope SAE features,
and show that a source-stable structured assistant packaging feature set
selectively tracks and weakly intervenes on a Gemma-based reward judge's
LLM-favoring margins. We then use deterministic text counterfactuals and
readout-space nulling to test whether those interpretable cue bundles drive
reward under approximate content preservation. The strongest current effect is
not generic formality but an assistant-style structure bundle; formal packaging
is smaller but specific, and benefit/value framing is not stable under matched
controls.
\end{quote}

This narrative keeps the current findings salient without overstating them.
The main open question for the final version is whether counterfactual audit
and readout nulling show enough cue-specific reward movement and preference
retention to be framed as a useful mitigation method rather than only as a
mechanistic diagnostic.

\bibliographystyle{plainnat}
\bibliography{references}

\end{document}
